% ============================================================
% data.tex
% ============================================================

This chapter describes the data sources, sample construction
procedure, and variable definitions underlying our analysis.
We document known data quality challenges in the Federal
Judicial Center's Integrated Database, explain how we address
them through composite identification strategies and
alternative coding schemes, and report the final sample
characteristics used in all subsequent analyses.

\section{Primary Data Source: FJC Integrated Database}

The Federal Judicial Center's Integrated Database (IDB)
aggregates case-level records from all 94 federal district
courts, reported by the Administrative Office of the
U.S.\ Courts \citep{fjc_idb_portal, fjc_idb_guide}. Civil
case records span from approximately 1970 to the present,
with substantially improved coverage quality after 1988.
The IDB is the most comprehensive source for population-level
data on federal civil litigation, covering virtually all cases
filed in federal district courts regardless of whether they
generate published opinions or disclosed settlement amounts.
This population-level coverage is a key advantage over
practitioner datasets, which typically track only high-profile
cases with publicly disclosed outcomes \citep{cornerstone2024}.

\subsection{Key Fields}

For each civil case, the IDB provides: filing date;
termination date (missing for pending cases, treated as
right-censored); nature-of-suit (NOS) code; disposition code;
district and circuit identifiers; origin code; class action
flag; procedural progress code; jurisdiction code; jury demand
indicator; monetary demand amount; and the \texttt{section}
field recording the statutory provision under which the case
was filed. Securities class actions are identified using
NOS Code~850 (``Securities, Commodities, and Exchange'').

\subsection{Data Quality and Known Limitations}

Empirical researchers have documented several data quality
issues directly relevant to our analysis
\citep{clermont2002, hadfield2004, scales2024}:

\begin{itemize}
    \item \textbf{Class action flag unreliability.} The IDB
    class action flag misses substantial fractions of actual
    class actions and incorrectly flags some non-class
    actions \citep{clermont2002, alexander1991}. Because the
    IDB case title field exhibits over 99\% missingness in our
    NOS~850 sample, keyword-based identification is not
    feasible with IDB data alone. CourtListener integration
    for title-based validation is deferred to future work
    (Chapter~\ref{ch:future}).

    \item \textbf{Disposition code ambiguity.} The IDB
    voluntary dismissal code (Code~12) is known to include
    hidden settlements: parties sometimes stipulate to
    voluntary dismissal as the formal docket entry while
    privately settling \citep{scales2024}. We address this
    through three alternative coding schemes described in
    Section~\ref{sec:coding_schemes}.

    \item \textbf{Missingness in secondary fields.} Monetary
    demand amounts exhibit 98.4\% missingness; jury demand
    information is missing for 50.7\% of cases. Critical
    fields for duration analysis---filing date, termination
    date, disposition code, circuit, and origin---are
    well-populated.
\end{itemize}

\section{Sample Construction}

We construct the analysis sample through three stages.

\textbf{Stage 1: NOS~850 Cases.} We begin with all civil
cases coded as NOS~850 filed from 1990 onward, yielding
65,899 cases. The 1990 start date balances sample size
against data quality and provides a meaningful pre-PSLRA
comparison period.

\textbf{Stage 2: Class Action Identification.} We retain
cases with IDB class action flag equal to 1, yielding
13,708 cases (20.8\% of NOS~850).

\textbf{Stage 3: Valid Duration.} We exclude cases with
non-positive computed durations (740 zero-duration
administrative artifacts, not substantive resolutions),
yielding a final sample of \textbf{12,968} federal securities
class actions filed between 1990 and 2024.

\begin{table}[htbp]
\centering
\caption{Sample Construction}
\label{tab:sample_construction}
\renewcommand{\arraystretch}{1.2} % Vertical breathing room
\begin{tabular*}{\textwidth}{@{\extracolsep{\fill}} l r r @{}}
\toprule
\textbf{Stage} & \textbf{N} & \textbf{\%} \\
\midrule
All NOS 850 cases (1990--present) & 65,899 & 100.0 \\
Class action flag $= 1$            & 13,708 &  20.8 \\
Valid duration ($> 0$ days)       & 12,968 &  19.7 \\
\bottomrule
\end{tabular*}
\end{table}

\section{Disposition Coding Schemes}
\label{sec:coding_schemes}

We implement three coding schemes that span the plausible
range of settlement/dismissal classifications given the
known Code~12 ambiguity:

\begin{itemize}
    \item \textbf{Scheme~A (Primary).} Code~13 $\to$
    settlement; Codes~2, 3, 4, 12, 14, 15, 17, 18, 19, 20
    $\to$ dismissal. Code~6 (judgment on motion before
    trial) is disaggregated using the IDB \texttt{JUDGMENT}
    field: \texttt{JUDGMENT}~$= 1$ (plaintiff victory,
    701~cases) $\to$ settlement; \texttt{JUDGMENT}~$= 2$
    (defendant victory, 1,418~cases) $\to$ dismissal;
    \texttt{JUDGMENT}~$\in \{3, 4, -8, \text{NA}\}$
    (ambiguous or missing, 270~cases) $\to$ censored.
    Trial codes, transfer codes, remands, and reopened
    cases are treated as censored, as they do not represent
    terminal resolution. This is the most conservative
    treatment and is used for all main results.

    \item \textbf{Scheme~B (Liberal).} Identical to
    Scheme~A except Code~12 is reclassified as settlement,
    reflecting the hypothesis that voluntary dismissals
    predominantly represent private settlements
    \citep{scales2024}.

    \item \textbf{Scheme~C (Expanded).} Identical to
    Scheme~B with the additional reclassification of
    Code~5 (consent judgment) as settlement.
\end{itemize}

\paragraph{Code~6 Disaggregation.}
In the raw IDB extract, disposition Code~6 (``Judgment on
Motion Before Trial'') encompasses 2,389 cases---18.4\% of
the analysis sample. Prior studies using the IDB typically
classify all Code~6 cases as dismissals, reasoning that
most judgments on pre-trial motions in securities
litigation are defense-favorable (granted motions to
dismiss or summary judgment). However, this default
classification conflates two substantively different
outcomes: plaintiff-favorable pre-trial judgments and
defendant victories. We resolve this measurement error by using the
IDB's \texttt{JUDGMENT} field, which records whether the
judgment favored the plaintiff (\texttt{JUDGMENT}~$= 1$) or
the defendant (\texttt{JUDGMENT}~$= 2$). This
disaggregation moves 701 previously hidden settlements into
their correct settlement pathway and retains 1,418
defendant victories as dismissals. The 270 cases with
ambiguous or missing \texttt{JUDGMENT} values are
conservatively censored. This refinement is applied
identically across all three coding schemes and addresses
the well-documented ``hidden settlements'' and ``hidden
dismissals'' measurement error in standard IDB extracts
\citep{scales2024}. We note that 93.4\% of the
reclassified plaintiff victories (655 of 701) are
post-PSLRA filings, which reflects the era-specific growth
in Code~6 usage rather than a systematic coding artifact;
the analytical implications of this temporal concentration
are discussed in Chapter~\ref{ch:discussion}.

\begin{table}[htbp]
\centering
\caption{Outcome Distribution by Disposition Coding Scheme}
\label{tab:schemes}
\renewcommand{\arraystretch}{1.2}
\begin{tabular*}{\textwidth}{@{\extracolsep{\fill}} l c c c c c c @{}}
\toprule
\textbf{Scheme} & \textbf{N} & \textbf{Settl.} & \textbf{Dism.} & \textbf{Cens.} & \textbf{\% S} & \textbf{\% D} \\
\midrule
A (Primary)  & 12,968 & 2,716 & 8,733 & 1,519 & 20.9 & 67.3 \\
B (Liberal)  & 12,968 & 4,520 & 6,929 & 1,519 & 34.9 & 53.4 \\
C (Expanded) & 12,968 & 4,834 & 6,929 & 1,205 & 37.3 & 53.4 \\
\bottomrule
\end{tabular*}
\end{table}

Under Scheme~A, 20.9\% of cases settle and 67.3\% are
dismissed, with 11.7\% censored. The Scheme~A settlement
rate is below the roughly 46\% reported by
Cornerstone Research \citep{cornerstone2024}, reflecting two
factors: Cornerstone tracks cases with publicly disclosed
settlements, while the IDB captures the full population
including small cases dismissed early; and Scheme~A
conservatively treats Code~12 as dismissal. Under Scheme~B,
which reclassifies Code~12 voluntary dismissals as
settlements, the settlement rate rises to 34.9\%, closing
much of the gap with industry benchmarks.

\section{Variable Definitions}

\subsection{Outcome Variable Construction}

The outcome variable for all survival models is the pair
$(T_i, \delta_i)$, where $T_i$ is the duration in years from filing 
to termination, and $\delta_i \in \{0, 1, 2\}$ is the event indicator 
under the applicable coding scheme. Cases still pending at the IDB 
snapshot are treated as right-censored ($\delta_i = 0$).

\subsection{PSLRA Regime}

The PSLRA indicator is a binary variable equal to 1 if the case 
\texttt{filedate} $\geq$ December~22, 1995, representing cases filed 
subject to the Act's heightened pleading standards.

\subsection{Circuit}

Circuit is coded from the IDB district identifier and treated
as a categorical variable. The Second Circuit serves as the
reference category in all regression models.

\subsection{Case Origin and MDL Flag}

We capture procedural pathways using two variables. First, the IDB 
origin code is categorized into: ``Original'' (federal filing, reference category), 
``Removed'' (from state court), ``MDL'' (transfers), and ``Other''. 
Second, because MDL cases can originate differently but share 
post-filing consolidation mechanics, we define a binary MDL flag 
equal to 1 if the \texttt{mdldock} field is populated with a 
non-missing, non-sentinel value. 

\subsection{Jurisdictional and Statutory Basis}

We construct a binary indicator for federal question jurisdiction 
using the IDB \texttt{juris} field. For statutory basis, we use the 
IDB \texttt{section} field to classify claims as: \texttt{0078} 
(Exchange Act \S~10(b)) $\to$ ``10(b)''; \texttt{0077} (Securities Act
\S~11) $\to$ ``Section~11''; and all other values $\to$ ``Other/Both''. 
A missingness indicator is included in all models to preserve the 
estimation sample.

\section{Sample Summary Statistics}

With our sample and variables defined, we present the baseline 
descriptive statistics that motivate our survival models. 

\subsection{Geographic Distribution}
The sample is heavily concentrated geographically. The Ninth Circuit 
accounts for 31.2\% (4,051 cases) and the Second Circuit accounts for 
28.8\% (3,737 cases), reflecting California and New York as the 
undisputed centers of U.S. securities litigation. Table~\ref{tab:bycircuit} 
details the full distribution.

\begin{table}[htb]
\centering
\caption{Sample Distribution by Circuit}
\label{tab:bycircuit}
\renewcommand{\arraystretch}{1.2}
\begin{tabular*}{\textwidth}{@{\extracolsep{\fill}} l r r @{}}
\toprule
\textbf{Circuit} & \textbf{N} & \textbf{\%} \\
\midrule
Ninth (CA)         & 4,051 & 31.2 \\
Second (NYC)       & 3,737 & 28.8 \\
Third (PA/NJ/DE)   &   828 &  6.4 \\
Fifth (TX/LA)      &   752 &  5.8 \\
Fourth             &   695 &  5.4 \\
Eleventh           &   627 &  4.8 \\
Seventh            &   535 &  4.1 \\
Tenth              &   532 &  4.1 \\
Eighth             &   453 &  3.5 \\
First              &   381 &  2.9 \\
Sixth              &   361 &  2.8 \\
D.C.               &    16 &  0.1 \\
\midrule
\textbf{Total}     & \textbf{12,968} & \textbf{100.0} \\
\bottomrule
\end{tabular*}
\end{table}

\subsection{Duration by Outcome}
Table~\ref{tab:duration} presents duration statistics for resolved 
cases under the primary conservative coding assumption (Scheme A). 
Settlements take substantially longer to resolve than
dismissals: median 2.81 years versus 1.40 years. The
dismissal distribution is strongly right-skewed, reflecting
a large mass of cases dismissed quickly on pleadings within
the first year, alongside a long right tail of cases
that survive initial motions.

\begin{table}[htbp]
\centering
\caption{Duration Statistics by Outcome (Resolved Cases, Scheme A)}
\label{tab:duration}
\renewcommand{\arraystretch}{1.2}
\begin{tabular*}{\textwidth}{@{\extracolsep{\fill}} l r r r r r @{}}
\toprule
\textbf{Outcome} & \textbf{N} & \textbf{Mean} & \textbf{Median} & \textbf{Q25} & \textbf{Q75} \\
\midrule
Settlement & 2,716 & 3.18 & 2.81 & 1.62 & 4.23 \\
Dismissal  & 8,733 & 2.02 & 1.40 & 0.29 & 3.00 \\
\bottomrule
\end{tabular*}
\end{table}



\subsection{Outcomes by PSLRA Regime}
Table~\ref{tab:byregime} illustrates the stark compositional shift 
in outcomes associated with the PSLRA. Of the 12,968 total cases,
1,032 (8.0\%) were filed pre-PSLRA and 11,936 (92.0\%) post-PSLRA.
The pre-PSLRA settlement rate sits at 40.0\%, whereas the post-PSLRA
settlement rate falls to 19.3\%. This raw disparity underscores 
the necessity of our formal hazard modeling in Chapter 5 to control 
for competing risks and concurrent case characteristics.

\begin{table}[htbp]
\centering
\caption{Outcome Distribution by PSLRA Regime (Scheme A)}
\label{tab:byregime}
\renewcommand{\arraystretch}{1.2}
\begin{tabular*}{\textwidth}{@{\extracolsep{\fill}} l r r r r r r @{}}
\toprule
\textbf{Regime} & \textbf{N} & \textbf{Settl.} & \textbf{Dism.} & \textbf{Cens.} & \textbf{\% S} & \textbf{\% D} \\
\midrule
Pre-PSLRA   & 1,032  & 413   & 510   & 109   & 40.0 & 49.4 \\
Post-PSLRA  & 11,936 & 2,303 & 8,223 & 1,410 & 19.3 & 68.9 \\
\midrule
\textbf{Total} & \textbf{12,968} & \textbf{2,716} & \textbf{8,733} & \textbf{1,519} & \textbf{20.9} & \textbf{67.3} \\
\bottomrule
\end{tabular*}
\end{table}


\section{Data Access and Replication}

The IDB is publicly available from the Federal Judicial Center
upon request at no cost. All data cleaning, variable
construction, and coding scheme implementation are documented
in the R analysis script and will be released as a public
replication package alongside the final thesis submission.
CourtListener data integration---which would enable
title-based class action identification, docket activity
counts, and attorney identifiers---is identified as a
priority extension in Chapter~\ref{ch:future}.